\raggedbottom
\definecolor{seqNineBlue}{HTML}{356A93}
\definecolor{seqNineOrange}{HTML}{C97024}
\definecolor{seqNineGreen}{HTML}{477A45}
\tikzset{
  seq nine cell/.style={draw=black!55,minimum width=.66cm,
    minimum height=.48cm,inner sep=1pt,font=\footnotesize},
  seq nine blue/.style={seq nine cell,draw=seqNineBlue,
    fill=seqNineBlue!12},
  seq nine orange/.style={seq nine cell,draw=seqNineOrange,
    fill=seqNineOrange!15,thick},
  seq nine green/.style={seq nine cell,draw=seqNineGreen,
    fill=seqNineGreen!13},
  seq nine flow/.style={-{Latex[length=1.7mm]},thick}
}
\lstdefinestyle{seq-nine-sql}{style=bioseq-cpp26,language=SQL,
  basicstyle=\small\ttfamily,columns=fullflexible,keepspaces=true}

\section{One sequence, several ways to read it}
\label{sec:sequence-abstractions}

Suppose a chromosome contains 100 million symbols, but an alignment step needs
only positions 500000 through 500099. Its input is the 100-symbol slice
\(S[500000,500100)\). Decoding the whole chromosome into one byte per symbol
would first create a 100 MB buffer for this small request. Repeating that
conversion for each candidate alignment region would waste memory and time.

The database introduces another difficulty: the same sequence can arrive as
a packed value, several stored fragments, or an expression such as a reversed
slice. The algorithms need the symbols in the right order in each case.
BioSeqDB separates a sequence's meaning from the way its symbols are stored.
The distinction has four parts.
\begin{itemize}
  \item \textbf{Value.} A sequence has an alphabet, a length, and an ordered
        list of symbols.
  \item \textbf{Storage.} Symbols may already be in memory or may need to be
        read from stored fragments.
  \item \textbf{View.} A view describes which symbols to read and their order.
        A slice can therefore describe part of a sequence without copying its
        symbols.
  \item \textbf{Consumer.} An algorithm states whether it can read a stream
        of symbols or needs one contiguous block of packed data, stored at
        adjacent memory addresses.
\end{itemize}

Here, \emph{lazy} means that the system delays reading, decoding, or copying
symbols until they are needed. Creating a lazy value can still allocate small
objects and retain access to its source. A lazy slice is useful because its
description is small even when its source is large. It does not promise that
every later read will avoid copying. Data is \emph{resident} when it is already
in memory. An \emph{owner} is the object responsible for keeping that memory
and any source resources alive until their last use.

\section{Packed symbols and a common C++ interface}
\label{sec:compact-sequence-range}
\label{sec:sequence-bit-layout}

An alphabet assigns each symbol an integer called its \emph{rank}. For DNA4,
the four ranks are \(A=0\), \(C=1\), \(G=2\), and \(T=3\). A rank is an
encoding of a symbol, not a score. Two bits are enough for these four values.
Other alphabets can require more bits and different complement rules.

BioSeqDB packs several ranks in an unsigned integer called a \emph{storage
unit}. For a unit of \(B\) bytes and \(b\) bits per symbol, it stores
\[
  L=\left\lfloor\frac{8B}{b}\right\rfloor
\]
symbols per unit. The first symbol uses the lowest \(b\) bits. Symbol position
\(i\) is in unit \(\lfloor i/L\rfloor\), at bit offset
\(b(i\bmod L)\). Reading its rank loads that unit, shifts it right by the
offset, and keeps the lowest \(b\) bits.

\begin{figure}[H]
\centering
\begin{tikzpicture}[x=1cm,y=1cm,font=\footnotesize]
  \node[anchor=east] at (2.1,2.85) {logical position};
  \foreach \i in {0,...,5}
    \node at (2.6+.76*\i,2.85) {\i};
  \node[anchor=east] at (2.1,2.2) {symbols};
  \foreach \i/\s in {0/A,1/C,2/G,3/T,4/G,5/A}
    \node[seq nine blue] at (2.6+.76*\i,2.2) {\texttt{\s}};
  \node[anchor=east] at (2.1,1.2) {packed ranks};
  \foreach \i/\r in {0/00,1/01,2/10,3/11,4/10,5/00}
    \node[seq nine blue,font=\ttfamily\footnotesize]
      at (2.6+.76*\i,1.2) {\r};
  \foreach \i in {6,7}
    \node[seq nine cell,dashed,fill=black!4]
      at (2.6+.76*\i,1.2) {--};
  \draw[thick] (2.20,.82) rectangle (5.26,1.58);
  \draw[thick] (5.26,.82) rectangle (8.32,1.58);
  \node at (3.73,.45) {unit 0: one byte};
  \node at (6.79,.45) {unit 1: one byte};
  \draw[seq nine flow] (2.28,-.08)--(5.13,-.08);
  \draw[seq nine flow] (5.40,-.08)--(8.20,-.08);
  \node at (5.25,-.48) {low bits to high bits within each unit};
  \node[anchor=west,align=left] at (9.0,2.1)
    {6 symbols\\2 storage units\\2 bytes};
  \node[anchor=west,align=left] at (9.0,.85)
    {2 empty symbol slots\\in the final unit};
\end{tikzpicture}
\caption{Packing \texttt{ACGTGA} with four DNA4 symbols per byte. The dashed
slots are padding, so they are not part of the sequence. The drawing orders
slots by logical position, from low bits to high bits.}
\label{fig:lsb-packed-layout}
\end{figure}

The C++ type \texttt{compact\_sequence\_range<R>} provides common operations
for several resident representations. Its compile-time properties, called
\emph{traits}, supply the alphabet, unit type, and packing rules. Algorithms
use those traits instead of assuming that
every input is DNA4 or that every symbol occupies one byte.

Three sizes must be kept separate.
\begin{itemize}
  \item \texttt{size()} counts storage units in a compact range.
  \item \texttt{bp\_size()} counts logical symbols, excluding padding.
        The name is retained for all supported sequence alphabets.
  \item \texttt{byte\_size()} counts bytes occupied by the units.
\end{itemize}
If the range contains \(U\) units and the last unit has \(P\) empty symbol
slots, then \(N=UL-P\) is its logical length. Figure~\ref{fig:lsb-packed-layout}
has \(U=2\), \(L=4\), and \(P=2\), hence \(N=6\). Unused high bits in
a unit are a separate issue: three-bit ranks in an eight-bit unit leave two
high bits that cannot hold another symbol. Unlike final padding, these bits
are unused in every unit.

\noindent\begin{minipage}{\linewidth}
The following example uses the actual range API at a consumer boundary. The
caller owns \texttt{output}, which has at least \(last-first\) bytes, and
passes bounds satisfying \(0\leq first\leq last\leq N\).
\begin{BioseqCppSnippet}[caption={A consumer requests only the ranks it needs.},
  label={lst:compact-range-interface}]
template<supports_compact_sequence_range R>
void read_fragment(const compact_sequence_range<R>& sequence,
                   std::size_t first, std::size_t last,
                   std::span<std::uint8_t> output) {
    // One rank per output byte, in forward order.
    sequence.template to_encoded_into<false>(output, first, last);
}
\end{BioseqCppSnippet}
\end{minipage}

For resident packed data, this call seeks to the first selected unit and
decodes \(last-first\) symbols. It does not first decode the preceding
symbols. The storage-unit width describes packing; coordinates, scores, and
hashes can use other integer widths.

The common interface also does not impose one ownership rule. A compact range
over a packed vector can own its payload. An adapter over a caller's alphabet
container or a validated PostgreSQL value borrows it. A borrowed range remains
valid only while that container or value is alive.
